\clearpage\thispagestyle{empty}
\pdfbookmark[0]{Quellen zum Pallas-Bericht vom November 1802}{pallas180211-sources}
\begin{center}\textbf{Quellen zum Pallas-Bericht vom November 1802}\end{center}
\noindent Diese KI-Übertragung in heutiges Deutsch folgt dem Auszug aus Franz Xaver von Zachs Pallas-Bericht vom November 1802 in \href{https://gdz.sub.uni-goettingen.de/id/PPN236007467}{Gauß’ \textit{Werken}, Band~VI} (Göttingen, 1874), S.~230--231. Der Auszug enthält von Zachs Erwiderung auf William Herschel und zitiert Gauß’ Brief vom 16. Oktober über die Einordnung von Ceres und Pallas und die Regel für die Abstände der Planeten.
\par\smallskip\noindent Der \href{https://api.digitale-sammlungen.de/iiif/image/v2/bsb10538598_00513/full/1450,/0/default.jpg}{Erstdruck} steht als Beitrag LII in der \textit{Monatlichen Correspondenz}, Band~VI, S.~499--505. Der Auszug reicht von der \href{https://api.digitale-sammlungen.de/iiif/image/v2/bsb10538598_00516/full/1450,/0/default.jpg}{Mitteilung über Herschels Sonderdruck auf S.~502} bis zum \href{https://api.digitale-sammlungen.de/iiif/image/v2/bsb10538598_00519/full/1450,/0/default.jpg}{Ende des Gauß-Zitats auf S.~505}. Der Sonderdruck ist Herschels \href{https://doi.org/10.1098/rstl.1802.0010}{\textit{Observations on the Two Lately Discovered Celestial Bodies}}. Auf der von Gauß angeführten \href{https://api.digitale-sammlungen.de/iiif/image/v2/bsb10538595_00604/full/1450,/0/default.jpg}{Seite vom Juni 1801} steht die numerische Regel, deren allgemeine Form Wurm zugeschrieben wird. Die Herausgeber der \textit{Werke} fügen den Verweis auf \href{https://api.digitale-sammlungen.de/iiif/image/v2/bsb10538319_00180/full/1450,/0/default.jpg}{Wurms Erörterung im \textit{Astronomischen Jahrbuch} für 1790, S.~168} hinzu.
\par\smallskip\noindent Eine frühere englische Wiedergabe eines Teils von Gauß’ Brief steht in Alexander von Humboldts \href{https://www.geology.19thcenturyscience.org/books/1858-Humboldt-Cosmos/Vol-IV/PDFpages/1858-Humboldt-Cosmos-IV-0118.pdf}{\textit{Cosmos}, Band~IV, S.~118}, übersetzt von E.~C. Otté und B.~H. Paul (\href{https://www.geology.19thcenturyscience.org/books/1858-Humboldt-Cosmos/Vol-IV/PDFpages/1858-Humboldt-Cosmos-IV-0002.pdf}{New York, Harper \& Brothers, 1858}). In diesem Auszug fehlt \(4+24\). Die vorliegende Übertragung gibt die ganze Folge und die ganze Passage nach den \textit{Werken} wieder. David C. Bossard erstellte die digitale Fassung von \textit{Cosmos} im Dezember 2005.
\par\smallskip\noindent Übertragung und Satz wurden mit KI erstellt. Die Digitalisate der Zeitschrift und des Jahrbuchs stammen von der Bayerischen Staatsbibliothek.
\clearpage
